% =====================================================================
\chapter{Exponenciales y logaritmos}
\label{tema:5}
% =====================================================================

El logaritmo convierte productos en sumas y potencias en productos. Esa propiedad lo hace omnipresente en el curso: la log-verosimilitud, la entropía, la perplejidad y la escala de las leyes de potencia son logaritmos. Este tema repasa las propiedades de exponenciales y logaritmos siguiendo el capítulo~6 de OpenStax~\cite{openstax}, introduce la función log-suma-exp siguiendo a Blanchard, Higham y Higham~\cite{blanchard} y fija las unidades de información (bits y nats) como en el capítulo~2 de Cover y Thomas~\cite{ct}. La construcción rigurosa de $e^x$ y $\ln x$ requiere el cálculo integral~\cite[cap.~18]{spivak}; aquí se admiten sus propiedades básicas.

\section{Potencias de exponente real y la función exponencial}

Para $a>0$, la potencia $a^r$ está definida para $r$ racional (\cref{tema:1}) y se extiende a todo $x\in\R$ de forma que la función $x\mapsto a^x$ es continua. Se admiten las siguientes propiedades, que extienden la \cref{prop:exponentes}.

\begin{proposicion}[Función exponencial]\label{prop:exp}
Sea $a>0$. Para todos $x,y\in\R$: $a^x>0$; $a^{x+y}=a^xa^y$; $(a^x)^y=a^{xy}$; $a^{-x}=1/a^x$. Si $a>1$, la función $x\mapsto a^x$ es estrictamente creciente; si $0<a<1$, es estrictamente decreciente. Si $a\neq1$, toma cada valor positivo exactamente una vez.
\end{proposicion}

El número $e=\lim_{n\to\infty}(1+\frac1n)^n\approx2{,}71828$ es la base natural; más en general, $e^x=\lim_{n\to\infty}(1+\frac xn)^n$ para todo $x$~\cite[cap.~18]{spivak}.

\begin{proposicion}\label{prop:ex-cota}
Para todo $x\in\R$, $e^x\ge1+x$.
\end{proposicion}

\section{Logaritmos}

\begin{definicion}[Logaritmo]\label{def:log}
Sean $a>0$, $a\neq1$ y $x>0$. El \emph{logaritmo en base $a$} de $x$ es el único $y\in\R$ tal que $a^y=x$; se escribe $y=\log_ax$. Se usan $\ln=\log_e$ (logaritmo natural), $\log_2$ (binario) y $\log_{10}$ (decimal).
\end{definicion}

Por definición, $a^{\log_ax}=x$ para $x>0$ y $\log_a(a^y)=y$ para todo $y$. El logaritmo solo está definido para argumentos positivos.

\begin{teorema}[Propiedades de los logaritmos]\label{teo:log}
Sean $a,b>0$ distintos de~1, $x,y>0$ y $r\in\R$. Entonces:
\begin{enumerate}
  \item $\log_a(xy)=\log_ax+\log_ay$ y $\log_a(x/y)=\log_ax-\log_ay$;
  \item $\log_a(x^r)=r\log_ax$;
  \item cambio de base: $\log_ax=\dfrac{\log_bx}{\log_ba}$; en particular, $\log_ab\cdot\log_ba=1$;
  \item si $a>1$, $x<y$ si y solo si $\log_ax<\log_ay$.
\end{enumerate}
\end{teorema}

\begin{corolario}\label{cor:lnx}
Para todo $x>0$, $\ln x\le x-1$, con igualdad si y solo si $x=1$. En consecuencia, $\ln x\ge1-\frac1x$.
\end{corolario}

\section{Ecuaciones exponenciales y logarítmicas}

Para resolver una ecuación exponencial se toman logaritmos o se cambia la incógnita ($t=a^x$); en las logarítmicas se usan las propiedades del \cref{teo:log} para llegar a $\log_aA=\log_aB$, que equivale a $A=B$, y se descartan las soluciones que hacen negativo o nulo algún argumento.

\addtocounter{definicion}{1}%

\addtocounter{definicion}{1}%

\section{La función log-suma-exp}

\begin{definicion}[Log-suma-exp]\label{def:lse}
Para $z=(z_1,\dots,z_n)\in\R^n$, $\displaystyle\LSE(z)=\ln\sum_{i=1}^ne^{z_i}$. La función \emph{softmax} es $\softmax(z)_i=e^{z_i}/\sum_je^{z_j}=e^{z_i-\LSE(z)}$.
\end{definicion}

\begin{proposicion}\label{prop:lse}
Para todo $z\in\R^n$ y $c\in\R$, con $M=\max_iz_i$:
\begin{enumerate}
  \item invariancia por traslación: $\LSE(z_1+c,\dots,z_n+c)=\LSE(z)+c$; en consecuencia, la softmax no cambia al sumar $c$ a todas las componentes;
  \item cotas: $M\le\LSE(z)\le M+\ln n$.
\end{enumerate}
\end{proposicion}

La \cref{prop:lse}(i) con $c=-M$ da el \emph{cálculo estable}: $\LSE(z)=M+\ln\sum_ie^{z_i-M}$. Todas las exponenciales son ahora $\le1$, de modo que no hay desbordamiento, y al menos una vale exactamente~1, de modo que el logaritmo no recibe un cero.

\addtocounter{definicion}{1}%

\section{Escalas logarítmicas, bits y nats}

Una cantidad que varía en muchos órdenes de magnitud se representa mejor por su logaritmo. Si $y=Cx^{k}$, entonces $\log y=\log C+k\log x$: en ejes doblemente logarítmicos una \emph{ley de potencia} es una recta de pendiente $k$, que es como se presentan las leyes de escala del módulo~16. Si $y=Ca^{x}$, entonces $\log y$ es lineal en $x$ (escala semilogarítmica).

\begin{definicion}[Unidades de información]
La \emph{información} de un suceso de probabilidad $p\in(0,1]$ es $-\log_2p$ \emph{bits} o $-\ln p$ \emph{nats}~\cite[cap.~2]{ct}. Por el cambio de base, $1\ \text{nat}=\frac1{\ln2}\ \text{bits}\approx1{,}4427$ bits.
\end{definicion}

\addtocounter{definicion}{1}%

\addtocounter{definicion}{1}%

\addtocounter{definicion}{1}%

